
\chapter{Interprétation de l'opérateur de \textit{dressing} en GHD}

En \textit{Generalized Hydrodynamics} (GHD), l’un des concepts centraux est
l’opérateur de \textit{dressing}. Cet opérateur traduit la manière dont une quantité
microscopique, définie sur une particule quasi-libre, est modifiée par la présence
des interactions à l’échelle macroscopique, en tenant compte de la structure des
quasi-particules et de leur densité d’occupation.

\subsection*{Définition}

Étant donné une fonction nue $h(\theta)$, définie sur l’espace des rapidités
$\theta$, son \textit{dressing} est défini par l’équation intégrale :
\begin{equation}
    h^{\mathrm{dr}}(\theta) = h(\theta) + \int \dd{\alpha}\,
    \varphi(\theta-\alpha)\, n(\alpha)\, h^{\mathrm{dr}}(\alpha),
\end{equation}
où :
\begin{itemize}
    \item $\varphi(\theta-\alpha)$ est le noyau de diffusion lié à la matrice de
    diffusion à deux corps du modèle intégrable ;
    \item $n(\alpha)$ est la fonction d’occupation (ou facteur de remplissage)
    caractérisant l’état macroscopique.
\end{itemize}

\subsection*{Interprétation physique}

L’opérateur de \textit{dressing} traduit le fait qu’une quasi-particule donnée
n’est jamais réellement libre dans un modèle intégrable : son énergie, sa
quantité de mouvement, ou encore sa vitesse de groupe sont modifiées par
l’ensemble des autres quasi-particules présentes dans l’état.
Ainsi :
\begin{itemize}
    \item $p^{\mathrm{dr}}(\theta)$ correspond à la quantité de mouvement
    effective d’une excitation ;
    \item $E^{\mathrm{dr}}(\theta)$ correspond à son énergie effective ;
    \item la vitesse de groupe se réécrit comme :
    \begin{equation}
        v^{\mathrm{eff}}(\theta) =
        \frac{(E')^{\mathrm{dr}}(\theta)}{(p')^{\mathrm{dr}}(\theta)}.
    \end{equation}
\end{itemize}

\subsection*{Résumé}

L’opérateur de \textit{dressing} joue donc un rôle analogue à une
\og renormalisation hydrodynamique \fg{} des observables : il encode les effets
collectifs des interactions entre quasi-particules et permet de relier les
quantités microscopiques aux grandeurs effectives pertinentes pour la dynamique
à grande échelle en GHD.

\section*{Interprétation de l’opérateur de \textit{dressing}}

En \gls{GHD}, l’opérateur de \textit{dressing} joue un rôle fondamental.

\paragraph{Idée physique.}
Dans un système intégrable, les quasi-particules ne sont jamais libres : elles
interagissent entre elles uniquement par des décalages de phase
(\emph{phase shifts}) lors des collisions.
Ainsi, la propagation d’une excitation dépend de la présence des autres
quasi-particules dans l’état macroscopique.

\paragraph{Quantités nues vs. quantités habillées.}
Chaque observable locale (énergie, impulsion, charges conservées, courants, etc.)
possède une valeur dite \emph{nue} (ou \emph{bare}).
L’opérateur de \textit{dressing} transforme cette quantité nue en une quantité
\emph{effective}, qui incorpore les effets collectifs des interactions avec toutes
les autres excitations présentes.

\paragraph{Conséquences.}
\begin{itemize}
    \item Le \textit{dressing} dépend de la distribution des quasi-particules
    (densités de rapidité) caractérisant l’état.
    \item Il modifie la vitesse de propagation des excitations, ce qui conduit à la
    notion de \emph{vitesse effective}.
    \item Il redéfinit aussi la valeur transportée des charges conservées et de leurs
    courants.
\end{itemize}

\paragraph{Résumé.}
L’opérateur de \textit{dressing} est la traduction mathématique des effets
d’interactions collectives entre quasi-particules.
Il permet de passer de la description en termes de particules « nues »
à celle des excitations « effectives », qui sont les véritables degrés de liberté
hydrodynamiques se propageant dans le fluide intégrable.

